\chapter{The Phonetic Alphabet, Morse Code, and the Q-Code}
\label{ch:phonetics-codes}

Three shared international conventions let amateurs of any nationality
communicate unambiguously, often across significant language barriers:
a standard way to spell letters aloud, a standard code for sending
text as sound, and a standard shorthand for common operating phrases.
All three are examined constantly by licensing exams worldwide, because
all three are genuinely universal.

\section{The Phonetic Alphabet}

Individual letters (especially in callsigns) are easily confused when
spoken plainly, particularly over a noisy or weak radio link, or between
speakers of different native languages. The \textbf{NATO phonetic
alphabet} (formally the International Radiotelephony Spelling Alphabet)
assigns each letter a distinct, unambiguous keyword chosen so that no
two words sound alike even under poor conditions.

\begin{center}
\renewcommand{\arraystretch}{1.05}
\begin{tabular}{@{}cl@{\hspace{2.2em}}cl@{\hspace{2.2em}}cl@{}}
\toprule
Ltr & Word & Ltr & Word & Ltr & Word \\
\midrule
A & Alfa     & J & Juliett   & S & Sierra \\
B & Bravo    & K & Kilo      & T & Tango \\
C & Charlie  & L & Lima      & U & Uniform \\
D & Delta    & M & Mike      & V & Victor \\
E & Echo     & N & November  & W & Whiskey \\
F & Foxtrot  & O & Oscar     & X & X-ray \\
G & Golf     & P & Papa      & Y & Yankee \\
H & Hotel    & Q & Quebec    & Z & Zulu \\
I & India    & R & Romeo     &   & \\
\bottomrule
\end{tabular}
\end{center}

\begin{examtip}
Note the standardized spellings \emph{Alfa} (not ``Alpha'') and
\emph{Juliett} (with a double T) --- deliberately chosen so the words
are pronounced consistently by speakers of many different languages,
none of whom necessarily share English spelling conventions. Exams
occasionally test this exact spelling detail.
\end{examtip}

Numbers are spoken as their normal names in most amateur usage
(``one,'' ``two,'' \dots, ``nine,'' ``zero''), though some formal
aviation/maritime contexts use modified pronunciations (``niner'' for
9) to avoid confusion with similar-sounding words; check your local
convention if operating in a formal net or emergency-communications
context.

\section{Morse Code (CW)}

\textbf{Morse code}, still widely used on amateur HF bands under the
name \textbf{CW} (continuous wave, referring to the on-off keyed
carrier used to send it, Chapter~\ref{ch:modulation}), represents each
letter, digit, and punctuation mark as a sequence of short
(\textbf{dits}) and long (\textbf{dahs}) elements, conventionally in a
3:1 duration ratio, with defined spacing between elements, letters, and
words.

\begin{center}
\renewcommand{\arraystretch}{1.05}
\begin{tabular}{@{}cc@{\hspace{1.5em}}cc@{\hspace{1.5em}}cc@{}}
\toprule
Ltr & Code & Ltr & Code & Ltr & Code \\
\midrule
A & .$-$ & J & .$---$ & S & ... \\
B & $-$... & K & $-$.$-$ & T & $-$ \\
C & $-$.$-$. & L & .$-$.. & U & ..$-$ \\
D & $-$.. & M & $--$ & V & ...$-$ \\
E & . & N & $-$. & W & .$--$ \\
F & ..$-$. & O & $---$ & X & $-$..$-$ \\
G & $--$. & P & .$--$. & Y & $-$.$--$ \\
H & .... & Q & $--$.$-$ & Z & $--$.. \\
I & .. & R & .$-$. & & \\
\bottomrule
\end{tabular}
\end{center}

Numerals follow a regular pattern: 1 is \texttt{.----}, 2 is
\texttt{..---}, counting down one dit at a time to 5 (\texttt{.....}),
then counting up one dah at a time to 0 (\texttt{-----}).

\begin{keyidea}[Timing, not just symbols]
Morse code proficiency is fundamentally about \emph{timing} recognition
--- a skilled operator hears whole letter and word \emph{sounds} (the
rhythm of ``dah-dit-dah-dit'' as ``C,'' not four separately counted
elements) rather than consciously decoding dits and dahs one at a time.
Standard spacing is one unit between elements of the same letter, three
units between letters, and seven units between words, all scaled by the
sending speed (measured in words per minute, WPM).
\end{keyidea}

\section{CW Operating Abbreviations}

To speed up telegraphy communication, standardized abbreviations
(many predating amateur radio, originating in commercial and military
telegraphy) shorten common words and phrases. A representative sample
of the most frequently used:

\begin{center}
\renewcommand{\arraystretch}{1.15}
\begin{tabularx}{\linewidth}{@{}l Y l Y@{}}
\toprule
Abbr. & Meaning & Abbr. & Meaning \\
\midrule
AGN & Again & OM & Old man (fellow amateur) \\
BTU & Back to you & PSE & Please \\
CUL & See you later & RIG & Radio equipment \\
FB & Fine business (good/great) & RST & Signal report (R/S/Tone) \\
GA & Go ahead / Good afternoon & TNX/TKS & Thanks \\
GE & Good evening & UR & Your / You're \\
GM & Good morning & WX & Weather \\
HI & Laughter (``hi hi'') & 73 & Best regards \\
HR & Here & 88 & Love and kisses (sparingly) \\
HW & How (copy)? & & \\
\bottomrule
\end{tabularx}
\end{center}

\section{The Q-Code}

The \textbf{Q-code}, standardized by the ITU, is a set of
three-letter groups (always beginning with Q) originally developed for
commercial Morse telegraphy to compress common phrases into a fixed,
language-independent shorthand. Each Q-code can function either as a
question or a statement depending on context (with or without a
question mark), and, while designed for CW, many have carried over into
everyday voice operating slang worldwide.

\begin{center}
\renewcommand{\arraystretch}{1.05}
\begin{tabular}{@{}lp{9.5cm}@{}}
\toprule
Code & Meaning \\
\midrule
QRG & What is my exact frequency? / Your exact frequency is \dots \\
QRL & Is this frequency in use? / This frequency is in use, please
do not interfere. \\
QRM & Are you being interfered with? / I am being interfered with
(man-made interference). \\
QRN & Are you troubled by static? / I am troubled by static
(natural/atmospheric noise). \\
QRO & Shall I increase power? / Increase power. \\
QRP & Shall I decrease power? / Decrease power (also used as a noun:
``QRP operating,'' meaning low-power operation). \\
QRT & Shall I stop transmitting? / Stop transmitting. \\
QRV & Are you ready? / I am ready. \\
QRX & When will you call again? / I will call you again at \dots \\
QRZ & Who is calling me? \\
QSB & Is my signal fading? / Your signal is fading. \\
QSL & Can you acknowledge receipt? / I acknowledge receipt (also
used as a noun: a ``QSL card'' confirms a contact). \\
QSO & Can you communicate with \dots directly? / I can communicate
with \dots (also used as a noun: ``a QSO'' means a contact/conversation). \\
QSY & Shall I change frequency? / Change frequency to \dots \\
QTH & What is your location? / My location is \dots \\
\bottomrule
\end{tabular}
\end{center}

\begin{examtip}
QTH, QSL, QRP, QSO, and QRM/QRN are used so pervasively --- even in
casual English-language voice conversation among amateurs worldwide ---
that they function almost as ordinary vocabulary within the hobby. Learn
these five thoroughly even before memorizing the rest of the list.
\end{examtip}

\section{The RST Signal Report}

Signal reports are conventionally given using the \textbf{RST} system:
\textbf{R}eadability (1--5, from unreadable to perfectly readable),
\textbf{S}trength (1--9, from faint to extremely strong), and
\textbf{T}one (1--9, CW only, describing the purity/quality of the
received tone; commonly omitted or given as ``9'' for voice modes,
where it does not apply in the same way). A report of ``59'' (voice) or
``599'' (CW) --- perfectly readable, very strong, perfectly pure tone
--- represents an excellent signal.

\begin{quickreview}
\item The NATO phonetic alphabet (Alfa, Bravo, Charlie, \dots) gives
unambiguous letter spelling across languages and conditions.
\item Morse code represents letters as dit/dah patterns in fixed
timing ratios; proficiency means recognizing rhythm, not counting
symbols.
\item Standardized CW abbreviations (FB, GA, TNX, 73, \dots) speed up
telegraphy and remain common in general amateur slang.
\item The Q-code (QTH, QSL, QRM, QSY, \dots) is an ITU-standardized,
language-independent operating shorthand usable on both CW and voice.
\item RST reports (Readability/Strength/Tone) give a standardized way
to describe received signal quality.
\end{quickreview}

\begin{practiceproblems}
\item Spell your own callsign (or a hypothetical one) using the NATO
phonetic alphabet.
\item Write the Morse code for the letters composing ``CQ.''
\item Explain the meaning of QRZ, QSY, and QRT.
\item A CW operator sends ``UR RST 599 HR, TNX FER CALL, 73.''
Translate this into plain English.
\item Explain why a language-independent code system like the Q-code is
particularly valuable in amateur radio, compared to using plain
national-language phrases.
\end{practiceproblems}
